% !TEX root =./ECE448_Practice_main.tex
%
% L6 practice set (AM Simulation). Short by design: the L6 Lab packet carries
% the hands-on work. Blank copy: \begin{solution}[h] reserves h of work
% space. SOLUTIONS copy (\iskey=1) prints the red worked solutions.

\fullwidth{\textbf{Learning Objective 3.1: } I can describe DSB-AM, DSB-SC, and SSB
and identify each on a spectrum.
\\[2pt]\textbf{Learning Objective 3.2: } I can implement an AM modulator and
demodulator in GNU Radio and verify demodulated audio against the source.
\\[4pt]\normalfont\small Show your work. Write a \textbf{Documentation} line at the
end (\textbf{None} if you did not collaborate).}

\begin{questions}

%%%%%%%%%%%%%%%%%%%%%%%%%%%%%%%%%%%%%%%%%
\question \textbf{Reading a spectrum.} A 1 kHz tone modulates a 100 kHz carrier.
\begin{parts}

	\part List the frequencies of the spectral lines for DSB-AM, DSB-SC, USB, and LSB.

	\begin{center}
	\renewcommand{\arraystretch}{1.25}
	\begin{tabular}{| l | p{1.5in} | p{1.5in} |}
	\hline
	Signal & Lines (kHz) & Bandwidth (kHz) \\\hline
	DSB-AM & & \\\hline
	DSB-SC & & \\\hline
	USB & & \\\hline
	LSB & & \\\hline
	\end{tabular}
	\end{center}

		\begin{solution}[0.8in]
		DSB-AM: 99, 100, and 101 kHz, 2 kHz wide. DSB-SC: 99 and 101 kHz, 2 kHz wide.
		USB: 101 kHz only, 1 kHz of message bandwidth. LSB: 99 kHz only, 1 kHz of
		message bandwidth.
		\end{solution}

	\part A spectrum display shows two lines of equal height at 14.997 and
	15.003 MHz and nothing between them. Name the modulation, the carrier frequency,
	and the message tone.
		\begin{solution}[0.9in]
		DSB-SC: the lines are symmetric about 15.000 MHz, so that is the carrier, and
		there is no carrier line. The tone is $15.003 - 15.000 = 3$ kHz.
		\end{solution}

	\part A waterfall shows a voice signal from 7.1970 to 7.1997 MHz that comes and
	goes with speech, with no steady line beside it. Voice runs from 300 Hz to 3 kHz.
	Name the modulation and the frequency a receiver must tune to.
		\begin{solution}[1.0in]
		SSB, because it is about 2.7 kHz wide, one-sided, and carrierless. The audio from
		300 Hz to 3 kHz lies \emph{below} the carrier, so it is LSB with the carrier at
		\eq{ 7.1970 + 0.0030 = 7.1997 + 0.0003 = 7.2000\ \text{MHz} }
		LSB below 10 MHz is the amateur convention.
		\end{solution}

\end{parts}

\newpage
%%%%%%%%%%%%%%%%%%%%%%%%%%%%%%%%%%%%%%%%%
\question \textbf{DSB-AM.} A transmitter sends
$s(t) = \left[1 + m\cos(2\pi f_m t)\right]\cos(2\pi f_c t)$, with $m$ set by a
slider.
\begin{parts}

	\part On the flowgraph's frequency sink you drag $m$ from 0.5 up to 1. Say which
	lines change and which do not.
		\begin{solution}[0.9in]
		The four sideband lines grow, because $m$ scales the message before it reaches
		the carrier. The carrier line at $f_c$ stays the same height, because it comes from
		the constant 1, which $m$ does not touch.
		\end{solution}

	\part The carrier is the tallest line on the display. Explain why it carries no
	information, and why AM broadcasting keeps it anyway.
		\begin{solution}[1.0in]
		The carrier line is the same whatever the message is, so it tells the listener
		nothing. It is there so the receiver can be an envelope detector, a diode and a
		capacitor, with no oscillator of its own to match.
		\end{solution}

	\begin{minipage}[t]{0.58\linewidth}
		\part A station carries audio up to 5 kHz. Find its bandwidth.
		\begin{solution}[0.6in]
		\eq{ 2W = 2 \times 5 = 10\ \text{kHz} }
		\end{solution}
	\end{minipage}\hfill%
	\begin{minipage}[t]{0.37\linewidth}
		\raggedleft \vspace{12pt}
		\ansbox[1.0in]{10 kHz}
	\end{minipage}

	\part The operator raises $m$ to 1.2. Explain in words what an envelope detector
	now outputs, and why it no longer matches the message.
		\begin{solution}[1.1in]
		At the message's troughs, $1 + 1.2\cos(2\pi f_m t)$ goes as low as $-0.2$. The
		carrier's amplitude cannot be negative, so the envelope is the absolute value,
		and the detector outputs $\vert 1 + m\cos(2\pi f_m t)\vert$: the troughs fold back
		up. The audio distorts, and the new frequencies spread the signal outside its
		channel.
		\end{solution}

\end{parts}

\newpage
%%%%%%%%%%%%%%%%%%%%%%%%%%%%%%%%%%%%%%%%%
\question \textbf{The envelope detector flowgraph.} The receiver in
\texttt{L06-dsb-am.grc} is Abs, then a low-pass filter, then Multiply Const. The
flowgraph runs at 48 kHz with a 10 kHz carrier and a message up to 2.5 kHz.
\begin{parts}

	\part The low-pass filter has a 3 kHz cutoff and a 1 kHz transition. Explain why
	that works: what must it pass, and what must it stop?
		\begin{solution}[1.0in]
		It must pass the envelope, which holds the message up to 2.5 kHz, and stop the
		higher-frequency content that Abs leaves behind. A 3 kHz cutoff with a 1 kHz
		transition passes everything up to about 2.5 kHz, which keeps both tones, and
		stops everything above about 3.5 kHz.
		\end{solution}

	\part Explain the Multiply Const of $\pi/2$ after the filter.
		\begin{solution}[0.8in]
		Averaging the rectified carrier leaves $2/\pi$ of its peak, so the filter's
		output is too small by that factor. Multiplying by $\pi/2$ restores the envelope's
		level so we can compare it with $1 + m\,x(t)$ on the time sink.
		\end{solution}

	\begin{minipage}[t]{0.58\linewidth}
		\part The time sink shows the detector output lagging the envelope by 57 samples.
		Find the lag in milliseconds, and name the block that causes it.
		\begin{solution}[1.0in]
		\eq{ \dfrac{57}{48000} = 1.19\ \text{ms} }
		The low-pass filter delays its output by half its length. The flowgraph's Delay
		block delays the reference by the same 57 samples so the traces line up.
		\end{solution}
	\end{minipage}\hfill%
	\begin{minipage}[t]{0.37\linewidth}
		\raggedleft \vspace{12pt}
		\ansbox[1.0in]{1.19 ms}
	\end{minipage}

	\part The flowgraph has an Audio Sink and no Throttle. Explain why.
		\begin{solution}[0.8in]
		The sound card takes exactly 48,000 samples per second, so the Audio Sink paces
		the flowgraph the way an SDR does. A Throttle would add a second clock competing
		with the sound card's and cause audio dropouts.
		\end{solution}

\end{parts}

\newpage
%%%%%%%%%%%%%%%%%%%%%%%%%%%%%%%%%%%%%%%%%
\question \textbf{Coherent and product detectors.}
\begin{parts}

	\part In \texttt{L06-dsb-sc.grc} you drag the LO phase slider from $0^\circ$ to
	$90^\circ$ with the audio on. Describe what the coherent receiver's trace on the
	time sink and the audio do.
		\begin{solution}[0.9in]
		The red trace shrinks toward zero as the phase error grows, and the tones get
		quieter. At $90^\circ$ the trace is flat at zero and the audio is silent. The pitch
		never changes; only the level does.
		\end{solution}

	\part Explain why a DSB-SC receiver on the air needs carrier recovery, and how FM
	stereo provides it for its DSB-SC subcarrier.
		\begin{solution}[1.0in]
		The receiver's LO must match the carrier in frequency and phase, and DSB-SC
		transmits no carrier to lock to; at $90^\circ$ the output vanishes. FM stereo sends
		a 19 kHz pilot tone, and the receiver doubles it to rebuild the 38 kHz LO in the
		right phase.
		\end{solution}

	\begin{minipage}[t]{0.58\linewidth}
		\part A USB signal carries a 1 kHz tone. The product detector's LO sits 150 Hz
		above the carrier. Find the output tone. Repeat for LSB.
		\begin{solution}[1.0in]
		The USB tone is at $f_c + 1000$; mixing with $f_c + 150$ leaves
		$1000 - 150 = 850$ Hz. The LSB tone is at $f_c - 1000$, which leaves
		$1000 + 150 = 1150$ Hz.
		\end{solution}
	\end{minipage}\hfill%
	\begin{minipage}[t]{0.37\linewidth}
		\raggedleft \vspace{12pt}
		USB: \ansbox[1.0in]{850 Hz}

		\vspace{4pt}
		LSB: \ansbox[1.0in]{1150 Hz}
	\end{minipage}

	\begin{minipage}[t]{0.58\linewidth}
		\part The SSB flowgraph's Delay block holds the message back
		$(129 - 1)/2 + 57$ samples before comparing it with the receiver's output. Find
		the delay and explain each term.
		\begin{solution}[1.1in]
		\eq{ 64 + 57 = 121\ \text{samples} }
		The 129-tap Hilbert filter delays the message by 64 samples, half its length,
		and the receiver's low-pass filter adds 57 more.
		\end{solution}
	\end{minipage}\hfill%
	\begin{minipage}[t]{0.37\linewidth}
		\raggedleft \vspace{12pt}
		\ansbox[1.0in]{121 samples}
	\end{minipage}

\end{parts}

\vspace{1.5em}
\noindent\textbf{Documentation:} \rule[-2pt]{0.6\linewidth}{0.4pt}

\end{questions}
